\section{Discussion}
\label{sec:discussion}
% DOMINIC WRITE DISCUSSION

From our evaluation of our comparison approach using PERMANOVA, we have found that using the Zhang-Shasha tree edit distance algorithm and our HMM distance cost function did not produce statistically significant results.
Due to our finding of a p-value of approximately \textbf{0.4}, we can conclude that our approach is likely due to random chance.
Further, the test statistic of \textbf{1.029482} means that there is effectively no difference in the within- and between-group distances.

There are likely several reasons why this is the case.
Most importantly, the flaw in using RNA clans as the category in the PERMANOVA test as previously mentioned.
Due to clans not being a purely structural classification, using it in the PERMANOVA test was likely meaningless.

Another potential reason is our small dataset.
We only considered 40 covariance models and 10 clans.
At the time, there are 4,170 RNA families and 147 clans on Rfam, meaning that there is a much larger dataset we could reach for.

Using all 4,170 RNA families (assuming each family on Rfam has a covariance model) would mean just under 8.7 million comparisons to make.
It is unlikely we could have completed our approach on the entire dataset in a reasonable time frame using the consumer hardware available to us.
Despite this limitation, it is possible that the PERMANOVA test on this larger dataset could produce statistically significant results.
However, due to the previously discussed issue with using clans as a category, this approach would likely still be flawed.

The second method of using correlation of RNA sequences comparison to covariance model comparison is a much more promising avenue for fairly evaluating our approach.
RNA sequence comparison or alignment is a well studied problem in bioinformatics, with a number of approaches and techniques.
Given enough time, we likely could have repurposed an existing tool for this.
We regret that we were unable to conduct this evaluation in time to include in this paper.

One benefit of our approach is the actual run time for comparing two models.
While we have no concrete data for any one particular run, being able to perform 780 comparisons in approximately 3 hours on consumer hardware represents a reasonable run time.
Due to our implementation being written in Python and utilizing no amount of parallelization or caching between comparisons, it is likely that there are significant performance gains to made.
As each comparison is independent of any other comparison, parallelization would be an easy optimization that would likely see time improvements, even on consumer hardware.

There are threats to our approach.
While Zhang-Shasha is a competitive algorithm for tree edit distance of ordered and labeled trees, it is possible that our chosen ordering of nodes from the covariance model negatively impacted our results.
It is also possible that the cost function we implemented to determine the cost of deleting, inserting, and re-labeling nodes was not optimal.
Our HMM fragment to HMM fragment is a probabilistic process involving produced sequence likely hood and could be potentially tuned for better or worse results.
Without a rigorous method of evaluation, these individual aspects of our approach are difficult to properly evaluate and potentially tune.

Overall, we find our approach to still have potentially despite the negative results from the PERMANOVA analysis.
Given more time and a better method of evaluation, we believe this approach could have more impactful results.
